\chapter{Introduction}\label{ch:introduction}

\chapterepigraph{Architecture begins where engineering ends.}{Walter Gropius}

\openingpara
  {The indoor environments in which people spend much of their lives}
  {are increasingly becoming interactive.\sidenotedown{3.15\baselineskip}{I use \emph{interactive indoor environments} to refer to indoor built environments in which spatial, environmental, and computational systems are entangled through sensing, automation, adaptation, feedback, or data-driven control.} In contemporary offices, universities, homes, and other public buildings, these environments sense occupancy, regulate temperature, adjust lighting, monitor air quality, automate access, and respond to patterns of use. Designed to support efficiency, comfort, sustainability, and adaptation, they are no longer only physical settings for everyday activity. They are becoming computationally mediated spaces that participate in how people work, move, rest, concentrate, socialise, and make sense of their surroundings.\sidenotedown{2.0\baselineskip}{The phrase \emph{computationally mediated} does not mean that the building is experienced primarily as a computer system. It points to the increasing presence of computation within the spatial and environmental conditions of everyday life.}}

The importance of these environments is twofold. First, people spend much of their everyday life indoors, which means that indoor environments have a persistent influence on bodily, affective, social, and cognitive experience. Second, as buildings become more interactive, they increasingly shape experience through forms of sensing, automation, feedback, and adaptation that are not always visible or directly understandable to occupants. An occupancy sensor that dims the lights, an interface that displays air quality, or an automated system that changes the environment without explanation may all become part of how a space is felt and interpreted. Interactive indoor environments therefore raise questions not only about technical performance, but also about how people encounter, make sense of, and inhabit spaces that respond around them.

\section{the problem}

The design of interactive indoor environments remains strongly shaped by optimisation.\sidenote{Optimisation is not treated here as a problem in itself. Buildings must support health, energy performance, accessibility, safety, and everyday functioning. The issue is that optimisation provides only a partial vocabulary for understanding lived experience.} Environmental qualities are often treated as variables to be measured, predicted, personalised, or controlled: temperature, lighting, sound, air quality, occupancy, and energy use. These concerns are necessary. Indoor environments need to support health, sustainability, accessibility, and everyday functioning. However, optimisation provides only a partial vocabulary for experience. It can describe acceptable thresholds and efficient operation, but it is less able to account for how a space becomes calm, exposed, lively, tiring, supportive, confusing, or indifferent.

This limitation matters because occupants do not simply receive environmental conditions. They sense them through their bodies, interpret them through spatial and social cues, and live with them through routines, adjustments, and choices. A workspace may meet thermal or acoustic expectations while still feeling tiring or misaligned with the activity at hand. A bright atrium may feel energising, exposed, or overwhelming depending on visibility, crowding, movement, and purpose. A quiet room may support concentration, but it may also make every chair scrape, cough, or whispered conversation more noticeable. These qualities shape how occupants decide where to sit, when to stay, when to leave, and how to behave.

Human Building Interaction, or HBI, provides the disciplinary context for this thesis because it examines the relations between people, built environments, and interactive technologies.\sidenote{Human Building Interaction is used here as the field that brings HCI into relation with architecture, building technologies, and everyday spatial experience. The thesis positions HBI not only as a domain of system design, but as a way to study how buildings are sensed, interpreted, and inhabited.} Interactive indoor environments are not simply technical systems to be controlled or interfaces to be used. They bring together spatial, environmental, social, and computational conditions that shape everyday experience over time. Interaction in these environments may involve explicit actions, such as using a display or adjusting a setting, but it may also emerge through movement, orientation, attention, proximity, occupancy, and the changing conditions of the space itself.\sidenote{This broad understanding of interaction is important for the thesis: interaction is not limited to direct input, screen-based control, or intentional system use. It may also occur through spatial behaviour, environmental response, and peripheral awareness.}

The problem this thesis addresses is how interactive indoor environments can be understood and designed as spaces of experience. This requires methods and concepts that can account for how occupants sense, interpret, and inhabit indoor environments through everyday practices. The challenge is to approach these environments not only as spaces that are measured and regulated, but also as spaces that are felt, interpreted, negotiated, and lived with over time.

\section{core thesis argument}

This dissertation argues that designing for experience in interactive indoor environments is less about optimising environmental conditions and more about enabling meaningful forms of encounter between occupants, spatial qualities, and computational systems. These encounters are formed through three related processes: sensing, interpreting, and inhabiting.\sidenote{\emph{Sensing}, \emph{interpreting}, and \emph{inhabiting} form the conceptual backbone of the thesis. They are not treated as separate stages, but as overlapping ways through which occupants encounter interactive indoor environments.}

Sensing refers to how occupants register spatial qualities such as light, sound, temperature, air, materiality, atmosphere, and temporality through sensory, bodily, and affective responses. Interpreting refers to how occupants make meaning from environmental cues, automation, system behaviour, bodily sensations, and social context. Inhabiting refers to how occupants live with indoor environments through choices, adjustments, workarounds, preferences, and patterns of use that develop over time.

Together, these three processes shift the focus of HBI. Instead of asking only whether an environment is comfortable, efficient, or adaptive, the thesis asks how occupants come to understand and live with interactive environments. Instead of treating occupants primarily as users of building systems or recipients of environmental control, it treats them as participants in the ongoing formation of indoor experience. The aim is not to abandon comfort, performance, or sustainability, but to expand the experiential scope through which interactive indoor environments are studied and designed.

The thesis develops this argument through two complementary strands of inquiry and a culminating design instantiation. The first strand examines occupant experience in interactive indoor environments, including through Affective Interaction as a methodological lens. The second investigates expert imaginaries of multisensory and biophilic interaction. The thesis then instantiates these empirical and conceptual insights through an exploratory AI supported sound system for shared indoor workspaces.

\subsection{understanding interactive indoor experiences through a new lens}

The first strand of the thesis draws on Affective Interaction as a methodological lens for studying indoor experience beyond conventional comfort assessment.\sidenote{Affective Interaction is used as a lens within one strand of the thesis, rather than as the organising framework for the whole dissertation. Its value here is methodological: it helps foreground affect as situated, relational, and interactional.} Affective Interaction is useful here because it does not treat affect as a fixed internal state to be detected or classified. Instead, it examines affect as situated, relational, and enacted through encounters between bodies, technologies, spaces, activities, and social situations.

This strand responds to a gap in how indoor environments are commonly evaluated. Comfort assessments and post-occupancy evaluations can reveal important information about satisfaction, preference, and perceived environmental quality. However, they often leave less room for ambiguity, temporality, atmosphere, social context, and the interpretive work through which occupants make sense of space.\sidenote{Post-occupancy and comfort-oriented approaches remain important, but they often privilege evaluation after use rather than the unfolding interpretation of space during everyday occupation.} Indoor experience is not always reducible to whether a condition is liked or disliked. It may involve uncertainty, adaptation, anticipation, comparison, tolerance, frustration, attachment, calm, or withdrawal.

Through empirical studies in interactive learning and working environments, this strand investigates how occupants describe their experiences, how they make sense of environmental and technological behaviour, and how affective qualities such as calm, exposure, privacy, frustration, and attentiveness emerge through situated use. These studies reveal fine-grained patterns of interaction and interpretation that are often overlooked by conventional comfort assessments. This account provides the empirical basis for rethinking what interactive indoor environments might be designed to support.

\subsection{imagining biophilic interaction}

The second strand of the thesis turns from understanding existing experience to imagining how interactive indoor environments could be designed otherwise. It investigates expert imaginaries of multisensory and biophilic interaction, with the aim of expanding the experiential vocabulary of technologically mediated indoor environments.\sidenote{I use \emph{imaginaries} to refer to future-oriented ways of reasoning about what interactive indoor environments could become. The term foregrounds design possibility rather than prediction.}

In this thesis, biophilic design is used as a generative lens for rethinking interaction in indoor environments through relations with nature and natural processes.\sidenote{Biophilic design is often reduced to visible natural elements such as plants, daylight, water, natural materials, or organic forms. This thesis uses it more broadly, as a way to reason about sensory variation, environmental process, atmosphere, and relations with natural phenomena.} Rather than treating biophilia as the addition of plants, daylight, water, natural materials, or organic forms, this strand asks how interactive environments might support qualities such as variation, ambiguity, refuge, prospect, movement, attention, immersion, distance, and connection.

Through expert interviews and design sessions, this strand examines how architects, designers, researchers, and practitioners imagine the future of multisensory indoor environments. These imaginaries reveal possibilities that are not adequately captured by dominant narratives of optimisation and automation. They point toward environments that are atmospheric rather than merely responsive, multisensory rather than visually or thermally dominated, and experiential rather than only functional.

The contribution of this strand is a biophilic interaction design framework for HBI. The framework structures expert insights into design opportunities for future interactive indoor environments, providing a way to move from empirical accounts of occupant experience toward design possibilities that are multisensory, situated, and environmentally responsive.

\subsection{instantiating the argument through sound}

The final study moves from understanding and imagining interactive indoor experience toward concrete design instantiation. The focus on sound follows from the preceding empirical and design-led work, where multisensory experience, atmosphere, ambience, and peripheral environmental qualities emerged as important but underdeveloped dimensions of interactive indoor environments.\sidenote{The focus on sound is not incidental. It follows from the thesis trajectory: the scoping review identifies acoustic and multisensory concerns within broader experience design, while the expert work foregrounds sound, atmosphere, and ambience as generative materials for biophilic interaction.} Sound is especially relevant because it is continuous, difficult to ignore, and deeply involved in how indoor spaces are felt.

Shared workspaces are shaped by HVAC, speech fragments, footsteps, typing, equipment noise, and moments of relative quiet. These sounds influence concentration, privacy, social awareness, fatigue, and atmosphere, yet they are often treated primarily as problems of noise, distraction, or acoustic control. This thesis instead examines sound as a material through which contextual and peripheral forms of interaction can be explored.\sidenote{Treating sound as a design material does not deny its role in disturbance or acoustic comfort. It expands the question from how sound can be reduced or controlled to how sonic qualities might participate in situated, multisensory experience.}

The thesis culminates in an exploratory AI supported sound system for shared indoor workspaces.\sidenote{The role of AI is design-oriented. In this prototype, AI is used to explore transformations between everyday indoor sound and nature-inspired ambient textures, rather than to classify occupants, infer emotional states, or optimise behaviour.} Developed as a grounded design instantiation of the preceding studies, the system explores how sonic features of everyday indoor environments can inform ambient audio textures inspired by natural soundscapes. Through an online survey and expert discussions, the study examines how such transformations may support contextual, peripheral, and non intrusive forms of interaction.

The aim is not to propose sound as a universal solution for indoor experience, but to demonstrate how the thesis’ approach can be implemented in a concrete system that treats indoor experience as multisensory and responsive to its environment.\sidenote{The sound system is positioned as a design instantiation rather than a final solution. Its purpose is to make the thesis argument testable and discussable through a concrete interactive system.}

\section{organization}

The remainder of this dissertation is organised as follows.

Chapter 2 reviews the relevant literature on Human Building Interaction, smart buildings, indoor environmental experience, Affective Interaction, and biophilic design. It establishes the theoretical and methodological context for studying interactive indoor environments as affective, multisensory, and interpretive spaces.

Chapter 3 presents the first empirical strand of the thesis. It identifies patterns of affective and interpretive participation in indoor environments and discusses their implications for HBI research and design.

Chapter 4 shifts from occupant experience to design imaginaries. It investigates expert perspectives on multisensory and biophilic interaction in indoor environments and examines how these imaginaries expand the experiential scope of smart buildings.

Chapter 5 synthesises these insights into a biophilic interaction design framework. The framework articulates design opportunities for shaping affective, multisensory, and temporally unfolding relations between occupants, environments, and interactive systems.

Chapter 6 presents the exploratory AI supported sound system for shared indoor workspaces. The chapter describes the motivation, design rationale, system development, online survey, and expert discussions, showing how sound can operate as a concrete design material for contextual and peripheral interaction in smart indoor environments.

Chapter 7 concludes the thesis by summarising the main contributions, discussing limitations, and outlining implications for future research in Human Building Interaction, Affective Interaction, biophilic design, and smart building practice.
